\chapter{Sensibilità alla soglia e verifiche aggiuntive}
\label{app:soglia}

Nel Capitolo~\ref{chap:risultati} un valore è segnalato quando lo score supera il 95-esimo percentile dei propri valori sugli anni di validazione. È una scelta convenzionale, non ottimizzata. Questa appendice verifica quanto i risultati dipendono dalla soglia: la Sezione~\ref{app:soglia-principali} ripete le analisi principali con il 99-esimo percentile, la Sezione~\ref{app:soglia-concordate} ricostruisce le segnalazioni concordate con il 97-esimo e il 99-esimo, e la Sezione~\ref{app:soglia-pr} varia la soglia di ogni metodo con il cubo concordato fisso. Le Sezioni~\ref{app:discriminatore} e~\ref{app:profilo-tipi} riportano due verifiche che completano il Capitolo~\ref{chap:risultati}: un esempio e la levigatezza per la componente del discriminatore, e il profilo di $G_{\mathrm{n}}$ per tipo di cluster.

\section{Le analisi principali con il 99-esimo percentile}
\label{app:soglia-principali}


Questa sezione ripete le analisi principali delle Sezioni~\ref{sec:ris-correlazione-spaziale}--\ref{sec:ris-confronti} con il 99-esimo percentile, cioè segnalando l'1\% dei valori degli anni di validazione, per verificare quanto le conclusioni dipendano dalla soglia (Tabella~\ref{tab:ris-soglia}). Le correlazioni fra G e D della Sezione~\ref{sec:ris-g-d-legame} usano tutti i valori e non dipendono dalla soglia, e non sono ripetute; non è ripetuta nemmeno l'analisi del dropout.

\begin{table}[htbp]
    \centering
    \footnotesize
    \caption{Le analisi principali con la soglia al 95-esimo e al 99-esimo percentile degli score sugli anni di validazione: percentuali per luogo e zone sugli anni di validazione, confronti e cluster sugli anni di test. La differenza relativa è quella della Sezione~\ref{sec:exp-misure}.}
    \label{tab:ris-soglia}
    \vspace{0.6em}
    \begin{tabular}{@{}L{0.22}L{0.40}rr@{}}
        \toprule
        Analisi & Misura & 95-esimo & 99-esimo \\
        \midrule
        Percentuale per fascia di latitudine e mese (Sez.~\ref{sec:ris-climatologia}) & G senza normalizzazione climatologica, intervallo fra fasce e mesi & 1.6--11.1\% & 0.2--3.6\% \\
         & $G_{\mathrm{n}}$, intervallo & 3.7--6.3\% & 0.5--2.1\% \\
         & differenza relativa, G $\rightarrow$ $G_{\mathrm{n}}$ & 0.39 $\rightarrow$ 0.11 & 0.70 $\rightarrow$ 0.28 \\
        \addlinespace
        Zone dei cluster (Sez.~\ref{sec:ris-g-d}) & cluster di $G_{\mathrm{n}}$ & 184877 & 70954 \\
         & $D_{\mathrm{reale}}$ nel nucleo, percentile mediano & 5 & 1 \\
         & $D_{\mathrm{reale}}$ nel contorno, percentile mediano & 12 & 4 \\
         & D nel nucleo, percentile mediano & 27 & 2 \\
        \addlinespace
        Confronti (Sez.~\ref{sec:ris-confronti}) & segnalazioni di $G_{\mathrm{n}}$ condivise con la variazione in tre ore & 53.9\% & 46.4\% \\
         & condivise con MCD & 52.1\% & 45.8\% \\
         & condivise con HBOS & 39.2\% & 24.3\% \\
         & segnalazioni esclusive di $G_{\mathrm{n}}$ & 18.1\% & 27.4\% \\
         & valori segnalati da almeno 5 metodi su 9 & 4.22\% & 0.72\% \\
        \addlinespace
        Cluster (Sez.~\ref{sec:ris-eventi}) & valori nei cluster con almeno 20 valori & 93\% & 85\% \\
         & durata dei cluster più grandi dei tre tipi principali & 23--40 giorni & 8--24 giorni \\
        \bottomrule
    \end{tabular}
\end{table}

\paragraph{La normalizzazione climatologica nella coda.} Con la soglia più alta la normalizzazione climatologica riduce ancora molto le differenze fra luoghi e stagioni, ma non le annulla: resta un eccesso di segnalazioni a sud, dove la percentuale media è circa il doppio di quella a nord: 1.4\% fra 20$^\circ$ e 30$^\circ$N contro 0.7\% fra 50$^\circ$ e 60$^\circ$N, con l'1\% atteso. La normalizzazione climatologica allinea media e deviazione standard di G in ogni cella, non la frequenza dei suoi valori più estremi; con il 5\% questa differenza si vede appena, con l'1\% sì.

\paragraph{Il comportamento del discriminatore.} Con la soglia più alta i cluster sono meno della metà, e il comportamento del discriminatore è lo stesso e più marcato. Nel nucleo $D_{\mathrm{reale}}$ scende dal percentile 5 all'1 e D dal 27 al 2. Anche nel contorno $D_{\mathrm{reale}}$ scende, da 12 a 4: con la soglia più alta i cluster sono più estremi, e anche le celle vicine sono fra i valori più alti di $G_{\mathrm{n}}$.

\paragraph{Cluster e catene.} Con la soglia più alta le catene si accorciano ma restano: i cluster più grandi dei tre tipi con più valori durano fra 8 e 24 giorni, contro 23--40 giorni con il 5\%, e i cluster con almeno 20 valori contengono ancora l'85\% dei valori segnalati, contro il 93\% con il 5\%. I tre tipi con più valori nei cluster di almeno 20 valori restano velocità verticale, precipitazione e copertura nuvolosa, ma la precipitazione passa al primo posto.

\paragraph{Confronti con le baseline.} Con la soglia più alta il quadro resta lo stesso, con un accordo più basso. La variazione in tre ore e MCD restano i metodi più vicini a $G_{\mathrm{n}}$, con il 46\% delle sue segnalazioni condivise; il metodo più lontano resta HBOS. I valori segnalati da almeno cinque metodi su nove sono lo 0.72\%: il calo è in gran parte automatico, perché ogni metodo segnala circa cinque volte meno. L'accordo resta molto sopra quello atteso per caso: con l'1\% due metodi indipendenti condividerebbero circa l'1\% delle segnalazioni, mentre $G_{\mathrm{n}}$ ne condivide il 46\% con la variazione in tre ore.

\paragraph{Sintesi.} Le conclusioni sul discriminatore non dipendono dalla soglia, e con una soglia più alta diventano più nette. Anche il quadro del confronto con le baseline non cambia, con un accordo più basso e una percentuale più alta di segnalazioni esclusive di $G_{\mathrm{n}}$. Le catene si accorciano, ma restano. L'effetto della normalizzazione climatologica dipende invece dalla soglia: la normalizzazione climatologica riduce le differenze fra luoghi e stagioni di circa tre quarti con il 5\% e di circa tre quinti con l'1\%.
% Numeri (5 novembre 2026, optuna-27). Quota per luogo e zone: validazione, outputs/comportamento_optuna27
% (i commenti delle Sezioni 6.3 e 6.4); quote per luogo da logs/luogo_g_grezzo_{5,1}pc.log e
% luogo_g_clim_{5,1}pc.log (48 caselle, una cifra decimale; differenza relativa = scarto medio assoluto
% dalla media delle caselle diviso per la media: 0.390 e 0.107 al 5%, 0.697 e 0.281 all'1%); zone da
% zone_g_dterms_5pc.csv e zone_g_dterms_1pc.csv (logs/zone_g_dterms_*pc.log). Confronti: test,
% outputs/confronti_good2/logs/comparison_5pc.log e comparison_1pc.log con MODEL optuna-27.
% Eventi più grandi dei tre tipi (regola del solo evento più grande, outputs/finale_optuna27/casi_5pc e
% casi_1pc, cases.csv): al 5% velocità verticale 555 ore (23.1 giorni), precipitazione 960 (40.0),
% copertura nuvolosa 603 (25.1); all'1% precipitazione 411 ore (17.1 giorni, 138 istanti), velocità
% verticale 186 (7.8, 63), copertura nuvolosa 564 (23.5, 189). Eventi all'1% (logs/casi_1pc.log): 317892
% eventi, 7435523 valori, 6323252 in eventi di almeno 20 valori (85.0%); valori negli eventi con più di 19
% valori: precipitazione 2353900, velocità verticale 1949717, copertura nuvolosa 504300, radiazione 353676.
% Al 5% gli eventi con almeno 20 valori contengono il 92.5% dei valori (33402993 su 36093783).
% Rapporto fra le quote medie a sud e a nord con G+clim all'1%: validazione 1.43 / 0.68 = 2.1 (sul test,
% outputs/finale_optuna27/logs/luogo_g_clim_1pc.log, 1.79 / 0.74 = 2.4).

\section{Le segnalazioni concordate al 97-esimo e al 99-esimo percentile}
\label{app:soglia-concordate}

Le segnalazioni concordate della Sezione~\ref{sec:ris-concordate} sono ricostruite, con le stesse regole, con la soglia di ogni metodo al 97-esimo e al 99-esimo percentile degli score sugli anni di validazione, cioè segnalando il 3\% e l'1\% dei valori; sono ripetuti anche l'inizio, la fine e lo spostamento lungo il moto della Sezione~\ref{sec:ris-concordate-inizio} (Tabella~\ref{tab:app-soglia-concordate}). La sovrapposizione, la probabilità di superamento e ciò che c'è attorno ai cluster non sono ripetuti.

\begin{table}[htbp]
    \centering
    \footnotesize
    \caption{Le segnalazioni concordate con la soglia di ogni metodo al 95-esimo, 97-esimo e 99-esimo percentile degli score sugli anni di validazione, sugli anni di test. Lo spostamento lungo il moto è positivo quando le segnalazioni stanno davanti al cluster concordato. ``Più avanti della baseline'' riporta l'intervallo fra le otto baseline.}
    \label{tab:app-soglia-concordate}
    \vspace{0.6em}
    \begin{tabular}{@{}L{0.58}rrr@{}}
        \toprule
        Misura & 95-esimo & 97-esimo & 99-esimo \\
        \midrule
        \multicolumn{4}{@{}l}{\emph{Cubo e cluster concordati}} \\
        Valori concordati [\%] & 6.20 & 3.54 & 1.05 \\
        Valori concordati senza la chiusura [\%] & 4.22 & 2.41 & 0.72 \\
        Voto decisivo di $G_{\mathrm{n}}$ [\% del cubo] & 8.4 & 9.2 & 11.5 \\
        Voto decisivo di PCA [\% del cubo] & 15.3 & 17.3 & 21.5 \\
        Cluster brevi & 27233 & 25392 & 16549 \\
        Cluster lunghi & 2142 & 1850 & 785 \\
        Valori concordati nei cluster lunghi [\%] & 85.4 & 74.5 & 50.6 \\
        Durata del cluster più lungo [giorni] & 97.8 & 97.4 & 40.5 \\
        \addlinespace
        \multicolumn{4}{@{}l}{\emph{$G_{\mathrm{n}}$ nei cluster brevi}} \\
        Spostamento lungo il moto, mediana [km] & $-4.4$ & $-3.9$ & $-2.8$ \\
        Cluster con $G_{\mathrm{n}}$ davanti [\%] & 46 & 46 & 47 \\
        Inizio: $G_{\mathrm{n}}$ prima / dopo MCD [\%] & 18 / 51 & 16 / 50 & 14 / 48 \\
        \addlinespace
        \multicolumn{4}{@{}l}{\emph{$G_{\mathrm{n}}$ nei cluster lunghi}} \\
        Spostamento lungo il moto, mediana [km] & 16.3 & 17.5 & 9.7 \\
        Cluster con $G_{\mathrm{n}}$ davanti [\%] & 62 & 64 & 61 \\
        Cluster con $G_{\mathrm{n}}$ più avanti della baseline [\%] & 57--65 & 56--67 & 55--68 \\
        Spostamento lungo il moto di MCD, mediana [km] & 8.8 & 9.0 & 1.5 \\
        Inizio: $G_{\mathrm{n}}$ prima / dopo MCD [\%] & 16 / 54 & 14 / 56 & 8 / 59 \\
        \bottomrule
    \end{tabular}
\end{table}

\paragraph{Il cubo segue la soglia.} Con soglie più alte il cubo concordato si riduce insieme alle segnalazioni di ogni metodo, dal 6.20\% dei valori al 3.54\% e all'1.05\%. Il voto decisivo di $G_{\mathrm{n}}$ cresce, dall'8.4\% all'11.5\% del cubo, ma resta molto più basso di quello di PCA, che passa dal 15.3\% al 21.5\%: a ogni soglia il riferimento dipende poco dal voto di $G_{\mathrm{n}}$.

\paragraph{La durata dei cluster lunghi dipende dalla soglia.} I cluster lunghi contengono l'85.4\% dei valori concordati con il 95-esimo percentile, il 74.5\% con il 97-esimo e il 50.6\% con il 99-esimo; il cluster più lungo passa da 97.8 a 40.5 giorni. Con soglie più alte le catene si accorciano: la durata dei cluster lunghi dipende in parte dalla soglia, non solo dai sistemi che li formano.

\paragraph{L'inizio.} A ogni soglia $G_{\mathrm{n}}$ non anticipa l'inizio dei cluster concordati. MCD comincia prima di $G_{\mathrm{n}}$ in circa metà dei cluster brevi e in più della metà dei lunghi; $G_{\mathrm{n}}$ comincia prima di MCD in meno di un cluster su cinque.

\paragraph{Lo spostamento lungo il moto.} Nei cluster lunghi $G_{\mathrm{n}}$ sta davanti a ogni soglia, di 16.3, 17.5 e 9.7~km in mediana, e sta più avanti di ciascuna baseline nel 55--68\% dei cluster. Con il 99-esimo percentile lo spostamento lungo il moto è più piccolo, anche per MCD, e i sistemi veloci sono pochi: solo 7 cluster si spostano di almeno due celle ogni tre ore. Nei cluster brevi $G_{\mathrm{n}}$ resta a ogni soglia leggermente dietro, di 2.8--4.4~km.

\paragraph{La fine.} La fine è il risultato che dipende di più dalla soglia. Nei cluster brevi, con il 95-esimo percentile $G_{\mathrm{n}}$ finisce dopo cinque baseline su otto; con il 97-esimo dopo quattro, HBOS, COPOD, Isolation Forest e la distanza dal normale, prima di tre e senza differenza significativa con PCA; con il 99-esimo prima di PCA, CBLOF, MCD e della variazione in tre ore, dopo COPOD e senza differenza significativa con HBOS e con la distanza dal normale. Nei cluster lunghi $G_{\mathrm{n}}$ finisce prima di sette baseline su otto con il 95-esimo percentile, di sei con il 97-esimo e di cinque con il 99-esimo; a ogni soglia finisce prima di PCA, CBLOF, MCD, della distanza dal normale e della variazione in tre ore.

\paragraph{Sintesi.} Tre risultati della Sezione~\ref{sec:ris-concordate} non dipendono dalla soglia: $G_{\mathrm{n}}$ non anticipa l'inizio dei cluster concordati, nei cluster lunghi sta davanti lungo il moto più di ogni baseline, nei brevi sta leggermente dietro. Dipendono invece dalla soglia la durata dei cluster lunghi e la fine nei cluster brevi: che $G_{\mathrm{n}}$ finisca dopo la maggior parte delle baseline vale con il 95-esimo percentile, non con il 99-esimo.
% Numeri: Tesi/segnalazioni_concordate/risultati.md §9 (cubo: summary.json e log del cubo; passo 3: log di
% check_era5_concordant_lead.py --kind analysed|long, 12 novembre 2026); server outputs/concordate_5pc_regola_tesi*,
% outputs/concordate_3pc_regola_tesi*, outputs/concordate_1pc_regola_tesi*. Fine, lunghi (test dei segni, p<0.01):
% 5% prima di tutte tranne COPOD; 3% prima di PCA, CBLOF, MCD, normale, variazione, HBOS (p 0.0097), iForest p 0.017
% non significativo, dopo COPOD; 1% prima di PCA, CBLOF, MCD, normale, variazione, dopo HBOS, COPOD, iForest.
% Classe di velocita' di almeno 2 celle ogni 3 h all'1%: 7 cluster.

\section{Curve precision--recall rispetto al cubo concordato}
\label{app:soglia-pr}

Con il cubo concordato fisso, quello del 95-esimo percentile, la soglia di ogni metodo varia dall'80-esimo al 99.9-esimo percentile del suo score sugli anni di validazione. A ogni soglia le segnalazioni del metodo sono chiuse come nel cubo, e se ne calcolano la precision e il recall (Sezione~\ref{sec:exp-baseline}) rispetto al cubo concordato e rispetto al cubo senza il voto del metodo, $C^{(-k)}$ (Sezione~\ref{sec:exp-concordate}), sugli anni di test (Figura~\ref{fig:app-pr}). Un metodo che segnalasse a caso avrebbe una precision pari alla percentuale di valori concordati, 0.062. Il cubo non è una verità: precision e recall misurano l'accordo con il cubo concordato, che per costruzione favorisce i metodi simili fra loro; il confronto senza il proprio voto toglie l'effetto del voto del metodo, non questa somiglianza.

\begin{figure}[htbp]
    \centering
    \phantomcaption
    \label{fig:app-pr}
    \begin{subfigure}[t]{0.480\textwidth}
        \centering
        \includegraphics[width=\linewidth]{era5_concordate_pr_a}
        \caption[Precision e recall di ogni metodo al variare della propria soglia, dall'80-esimo al 99.9-esimo percentile, rispetto al cubo concordato al 95-esimo percentile, sugli anni di test; il pallino indica il 95-esimo percentile.]{Precision e recall di ogni metodo al variare della propria soglia, dall'80-esimo al 99.9-esimo percentile, rispetto al cubo concordato al 95-esimo percentile, sugli anni di test; il pallino indica il 95-esimo percentile. Al 95-esimo $G_{\mathrm{n}}$ è penultimo.}
        \label{fig:app-pr-a}
    \end{subfigure}
    \hfill
    \begin{subfigure}[t]{0.480\textwidth}
        \centering
        \includegraphics[width=\linewidth]{era5_concordate_pr_b}
        \caption[Come la Figura~\ref{fig:app-pr-a}, rispetto al cubo concordato senza il voto del metodo.]{Come la Figura~\ref{fig:app-pr-a}, rispetto al cubo concordato senza il voto del metodo. Ai percentili più alti la precision di $G_{\mathrm{n}}$ supera quella di CBLOF, della distanza dal normale e della variazione in tre ore.}
        \label{fig:app-pr-b}
    \end{subfigure}
\end{figure}

\paragraph{Al 95-esimo percentile $G_{\mathrm{n}}$ è penultimo.} Alla soglia del cubo $G_{\mathrm{n}}$ ha precision 0.594 e recall 0.697, le più basse dopo quelle della variazione in tre ore; senza il proprio voto 0.522 e 0.669, con lo stesso ordine. È atteso: sei dei nove metodi leggono una cella in un istante e segnalano in gran parte gli stessi valori, mentre $G_{\mathrm{n}}$ si sovrappone anche alla variazione in tre ore (Sezione~\ref{sec:ris-confronti}).

\paragraph{Il proprio voto pesa meno per $G_{\mathrm{n}}$.} Senza il proprio voto la precision di $G_{\mathrm{n}}$ cala di 0.072, quella delle altre baseline fra 0.077 e 0.131; cala meno solo quella della variazione in tre ore, di 0.051. È coerente con il voto decisivo di $G_{\mathrm{n}}$, più basso di quello di ogni modello shallow (Sezione~\ref{sec:ris-concordate-cubo}).

\paragraph{Ai percentili più alti la precision di $G_{\mathrm{n}}$ sale.} Con soglie più alte la precision di $G_{\mathrm{n}}$ senza il proprio voto sale fino a 0.92 al 99.9-esimo percentile: i suoi valori più estremi sono quasi sempre segnalati anche da almeno cinque degli altri otto metodi. Dal 99-esimo percentile supera quella di CBLOF, della distanza dal normale e della variazione in tre ore, che al 99-esimo e al 99.5-esimo percentile restano fra 0.68 e 0.80 (la variazione in tre ore arriva a 0.89 al 99.9-esimo); PCA, HBOS, COPOD, Isolation Forest e MCD arrivano fra 0.91 e 1.00.

\paragraph{Ai percentili più alti il recall è limitato dal numero di segnalazioni.} Al 99-esimo percentile ogni metodo segnala fra l'1.3\% e l'1.7\% dei valori, contro il 6.20\% del cubo: il recall non può superare 0.21--0.28, secondo il metodo, e i metodi vi si equivalgono, fra 0.19 e 0.24.
% Numeri: Tesi/segnalazioni_concordate/risultati.md §12; dati figures/dati/era5_concordate_pr.csv (trascritti dal log
% del server outputs/concordate_5pc_regola_tesi_pr.log, 13 novembre 2026; codice check_era5_concordant_pr.py, commit
% 576547c). Al 95: precision / recall G_n 0.594 / 0.697, senza voto 0.522 / 0.669; calo della precision: PCA 0.131,
% iForest 0.105, HBOS 0.093, MCD 0.091, COPOD 0.090, CBLOF 0.081, normale 0.077, G_n 0.072, variazione 0.051.
% Senza voto, 99 / 99.5 / 99.9: G_n 0.819 / 0.877 / 0.918; CBLOF 0.767 / 0.762 / 0.792; normale 0.682 / 0.689 / 0.786;
% variazione 0.715 / 0.800 / 0.891; al 99 gli altri cinque 0.908-0.985. Al 99 segnalati 1.30-1.73%, recall 0.191-0.244.
% Figura: figures/genera_figure_concordate.py, pr_figure.

\section{Che cosa giudica il discriminatore: un esempio e la levigatezza}
\label{app:discriminatore}

Questa sezione completa la Sezione~\ref{sec:ris-g-d} con un istante letto su una mappa e con la verifica della levigatezza delle ricostruzioni.

\paragraph{Un esempio: l'istante con G più alto.} La Figura~\ref{fig:ris-d-termini-mappa} mostra G e D il 9 agosto 2003 alle 12~UTC, l'istante con la media di G più alta degli anni di validazione. Nell'area principale del Nord Africa in cui G è alto, $D_{\mathrm{reale}}$ è vicino a zero, con una mediana di 0.07, e anche $D_{\mathrm{generato}}$ è fra i valori più bassi dell'istante, con una mediana di 0.19: il discriminatore giudica reali sia l'osservazione sia la ricostruzione, l'osservazione di più, e D è fra i valori più bassi del dominio. Su un'area più piccola più a sud-ovest, anch'essa con G alto, il giudizio sull'osservazione è meno netto ($D_{\mathrm{reale}}$ ha una mediana di 0.15), ma anche lì D è bassa. Dentro la linea della Figura~\ref{fig:ris-d-termini-mappa}, che racchiude l'1\% delle celle con G più alto, il percentile mediano di D fra le celle dell'istante è 6. $D_{\mathrm{reale}}$ è basso anche nella penisola arabica, dove G è intorno al 94-esimo percentile dell'istante.

\begin{figure}[htbp]
    \centering
    \phantomcaption
    \label{fig:ris-d-termini-mappa}
    \begin{subfigure}[t]{0.493\textwidth}
        \centering
        \includegraphics[width=\linewidth]{era5_d_termini_mappa_a}
        \caption[G il 9 agosto 2003 alle 12~UTC, l'istante con la media di G più alta degli anni di validazione, come percentile fra le celle dell'istante; la linea arancione racchiude l'1\% delle celle con G più alto.]{G il 9 agosto 2003 alle 12~UTC, l'istante con la media di G più alta degli anni di validazione, come percentile fra le celle dell'istante; la linea arancione racchiude l'1\% delle celle con G più alto. Linee di costa da Natural Earth.}
        \label{fig:ris-d-termini-mappa-a}
    \end{subfigure}
    \hfill
    \begin{subfigure}[t]{0.467\textwidth}
        \centering
        \includegraphics[width=\linewidth]{era5_d_termini_mappa_b}
        \caption[D il 9 agosto 2003 alle 12~UTC, come percentile fra le celle dell'istante; la linea arancione racchiude l'1\% delle celle con G più alto.]{D il 9 agosto 2003 alle 12~UTC, come percentile fra le celle dell'istante; la linea arancione racchiude l'1\% delle celle con G più alto. Dentro la linea D è fra i valori più bassi. Linee di costa da Natural Earth.}
        \label{fig:ris-d-termini-mappa-b}
    \end{subfigure}
% Percentile di D fra le celle dell'istante, dentro l'1% di G più alto (figures/dati/era5_g_d.npz, riga 1, 10 novembre 2026): mediana 6.3, il 98% sotto 25.
\end{figure}

\paragraph{Le ricostruzioni non sono più lisce delle osservazioni.} Una spiegazione possibile del comportamento di D è che il discriminatore riconosca le ricostruzioni dal loro aspetto più liscio: i modelli addestrati con un errore quadratico producono campi troppo regolari, e le ricostruzioni di un autoencoder variazionale smussano gli estremi (Sezione~\ref{sec:rw-meteo}). La rugosità (Sezione~\ref{sec:exp-misure}) è misurata su 240 istanti degli anni di validazione (Tabella~\ref{tab:ris-levigatezza}) con il modello senza dropout, per non mescolarla con il rumore delle maschere di dropout. In mediana sugli istanti, la rugosità delle ricostruzioni è però 1.10 volte quella osservata: l'ipotesi non regge.
% Testo completo prima del taglio del 9 novembre 2026: \paragraph{Le ricostruzioni non sono più lisce delle osservazioni.} Una spiegazione possibile del comportamento di D è che il discriminatore riconosca le ricostruzioni dal loro aspetto più liscio: in letteratura i modelli addestrati con un errore quadratico~\cite{rasp2020weatherbench,lam2023graphcast}, come in parte il generatore della tesi, e le ricostruzioni di un autoencoder variazionale~\cite{pacal2026heatwavedrivers} producono campi troppo regolari e smussano gli estremi. L'ipotesi è messa alla prova su 240 istanti della validazione (Tabella~\ref{tab:ris-levigatezza}), misurando la rugosità, cioè la differenza quadratica media fra celle adiacenti; il modello è usato senza dropout, perché il giudizio non cambi con la maschera di dropout. Le ricostruzioni non sono però più lisce delle osservazioni: la loro rugosità è in mediana 1.10 volte quella osservata, ed è minore solo nel 22\% degli istanti. La levigatezza della ricostruzione non spiega quindi il comportamento di D.

\begin{table}[htbp]
    \centering
    \footnotesize
    \caption{Rugosità e giudizio medio del discriminatore sull'osservazione e sulla ricostruzione, intere e lisciate: mediane su 240 istanti degli anni di validazione, con il modello senza dropout.}
    \label{tab:ris-levigatezza}
    \vspace{0.6em}
    \begin{tabular}{@{}lrr@{}}
        \toprule
        Campo & Rugosità & \shortstack[r]{Giudizio medio\\(1 = generato)} \\
        \midrule
        Osservazione & 0.00082 & 0.502 \\
        Ricostruzione & 0.00089 & 0.502 \\
        Osservazione lisciata & 0.00021 & 0.512 \\
        Ricostruzione lisciata & 0.00019 & 0.513 \\
        \bottomrule
    \end{tabular}
\end{table}

\paragraph{Una lettura: le celle rugose sono giudicate reali.} Le celle con la rugosità locale più alta sono giudicate più reali (correlazione di Spearman mediana $-0.61$ fra rugosità locale e giudizio sull'osservazione) e tendono ad avere G alto (correlazione mediana $0.49$). Lisciare l'osservazione la fa giudicare appena più generata: su tutte le celle il giudizio medio passa, in mediana sugli istanti, da 0.502 a 0.512 (Tabella~\ref{tab:ris-levigatezza}), e sale nel 99.6\% dei valori con G nel 5\% più alto, sia sugli anni di validazione sia su quelli di test.
% Rugosità locale: scripts/check_era5_discriminator_smoothness.py, local_roughness (quadrato della distanza dalla media della patch 3x3, media sulle variabili). 0.502 -> 0.512 è su tutte le celle (tabella); il 99.6% è la quota delle celle del 5% più alto di G in cui il giudizio sale (testo completo del 9 novembre, sotto).
 La rugosità potrebbe quindi spiegare in parte perché, dove G è più alto, il discriminatore giudica reali le osservazioni. È una lettura di correlazioni, non una dimostrazione: il lisciamento riduce anche i valori più estremi, un effetto che la verifica non separa.
% Testo completo prima del taglio del 9 novembre 2026: \paragraph{Una lettura: i campi rugosi sono giudicati reali.} La stessa verifica suggerisce un'altra lettura. Lisciando l'osservazione con la media della sua patch $3\times3$, il discriminatore la giudica più generata nel 99\% delle celle, e nel 99.6\% di quelle con G più alto, anche se di poco: in media il giudizio passa da 0.502 a 0.512. La rugosità locale ordina il giudizio sulle osservazioni: la correlazione di Spearman mediana fra le due è $-0.61$, negativa in tutti gli istanti, cioè le celle più rugose sono giudicate più reali. Le stesse celle tendono ad avere anche G alto, con una correlazione mediana di $0.49$: dove il campo è rugoso, il generatore tende a sbagliare di più. Coerentemente, nei 240 istanti il giudizio sull'osservazione vale in media 0.45 nel 5\% più alto di G, contro 0.50 su tutte le celle. La rugosità potrebbe quindi spiegare in parte perché, dove G segnala, il discriminatore giudica reali le osservazioni. È una lettura di correlazioni, non una dimostrazione: il lisciamento, per esempio, riduce anche i valori più estremi, un effetto che la verifica non separa.
% Numeri dell'esempio e della levigatezza:
% Mappa del 9/8/2003: mappe_istanti.npz e outputs/finale_optuna27/logs/mappe_istanti.log
% (G medio più alto = 9/8/2003 12 UTC, 0.00746); celle nell'1% più alto di G dell'istante
% (la linea arancione della figura): 85 celle 27.5-30.5N 2.5-12.5E, mediane D(reale) 0.073,
% D(generato) 0.186 (percentile 1 dell'istante), D -0.112 (percentile 7); 22 celle
% 23-25.5N 5.5-2W, D(reale) 0.151, D(generato) 0.265, D -0.116. Penisola arabica (lat <=
% 38N, lon >= 35E, D(reale) < 0.3): 164 celle, G al percentile 93.7 dell'istante
% (mediana), D(reale) 0.240; il minimo di D dell'istante è a 32.5N 39.5E.
% Levigatezza: scripts/check_era5_discriminator_smoothness.py, logs/levigatezza_d.log,
% levigatezza_d.csv (240 istanti, uno ogni 61, modo eval): rugosità mediana osservato
% 0.00082, ricostruito 0.00089, osservato lisciato 0.00021, ricostruito lisciato 0.00019;
% rapporto ricostruito/osservato mediana 1.101, sotto 1 nel 21.7% degli istanti; D medio
% 0.5020, 0.5019, 0.5118, 0.5134; D osservato nel 5% di G più alto 0.4549. Lisciando: D
% sale nel 99.3% delle celle (+0.0098), nel 5% di G più alto nel 99.6% (+0.0379).
% Spearman rugosità locale - D osservato -0.613 (negativa nel 100% degli istanti), - G
% +0.485, - aumento di D +0.762.
% Paçal et al., p. 971: 'the VAE reconstructions tend to smooth out extremes'.
% Levigatezza sul test (logs/levigatezza_d.log, 240 istanti, uno ogni 255, modo eval):
% rugosità mediana osservato 0.00087, ricostruito 0.00095, osservato lisciato 0.00022,
% ricostruito lisciato 0.00019; rapporto mediana 1.077, sotto 1 nel 31.7% degli istanti;
% D medio 0.5005, 0.5010, 0.5113, 0.5131; D osservato nel 5% di G più alto 0.4500.
% Lisciando: D sale nel 99.3% delle celle (+0.0105), nel 5% di G più alto nel 99.6%
% (+0.0432). Spearman rugosità locale - D osservato -0.611 (negativa nel 100% degli
% istanti), - G +0.485, - aumento di D +0.757.

\section{Il profilo di $G_{\mathrm{n}}$ per tipo di cluster}
\label{app:profilo-tipi}

Questa sezione completa la Sezione~\ref{sec:ris-eventi} leggendo il profilo dei cluster di $G_{\mathrm{n}}$ per tipo di cluster.

\paragraph{Per tipo di cluster, $G_{\mathrm{n}}$ cala in modo simile, come diverse baseline.} Vicino al bordo $G_{\mathrm{n}}$ è quasi uguale per tutti i tipi: una cella fuori vale fra 89 e 91. Più lontano si separa di poco: a sei celle va da 43 per la radiazione solare a 58 per la pressione. Non è una proprietà del solo $G_{\mathrm{n}}$. A due celle fuori dal bordo dei propri cluster (Figura~\ref{fig:ris-eventi-profilo-tipi-c}, senza i tipi con meno di 1000 valori a quella distanza), il percentile mediano di $G_{\mathrm{n}}$ va dal 74 all'82 a seconda del tipo. Intervalli altrettanto stretti hanno la variazione in tre ore (da 60 a 67), HBOS (da 73 a 81) e COPOD (da 76 a 85); la distanza dal normale, MCD e PCA vanno invece da poco più di 60 a oltre 80. Per sette metodi su nove il tipo più basso è la radiazione solare.
% Figura 6.13c: genera_figure_eventi.py, types_spread_figure, dai profili per tipo di ogni metodo sui propri cluster (figures/dati/baseline/profilo_*_5pc.csv ed era5_eventi_profilo.csv), posizione 2, tipi con almeno 1000 valori: G_n 74-82, variazione in tre ore 60-67, HBOS 73-81, COPOD 76-85, Isolation Forest 70-83, CBLOF 65-81, distanza dal normale 64-84, MCD 61-82, PCA 62-84.

\paragraph{Lo scarto standardizzato della variabile del tipo cala invece in modo diverso.} Le differenze grandi riguardano lo scarto standardizzato della variabile del tipo (Figura~\ref{fig:ris-eventi-profilo-tipi}). Nella figura sono in colore i tre tipi con più valori nei cluster di almeno 20 valori, velocità verticale, precipitazione e copertura nuvolosa, e la temperatura a due metri; ogni punto è la mediana dei valori a quella distanza in tutti i cluster del tipo, e i punti con meno di 1000 valori non sono tracciati, perché la mediana sarebbe poco affidabile. A differenza della distanza dal normale, che riassume tutte le variabili di una cella, lo scarto standardizzato riguarda una sola variabile. La variabile del tipo è per definizione la più lontana dal normale nel nucleo: è in parte atteso che lì lo scarto standardizzato sia alto e fuori scenda, e il confronto utile è fra i tipi. Il livello di fondo (Sezione~\ref{sec:exp-misure}) non è mai zero: vale per esempio 0.42 per la precipitazione e 0.98 per la pressione. Per precipitazione, copertura nuvolosa e radiazione solare lo scarto standardizzato torna vicino al livello di fondo entro due celle dal bordo: per la precipitazione scende da 1.33 nell'ultima cella del nucleo a 0.48 due celle fuori, per la copertura nuvolosa da 1.62 a 0.52. La velocità verticale scende più lentamente, da 1.38 a 0.78. Per la temperatura a due metri resta alto a lungo: a sette celle fuori lo scarto standardizzato è ancora 1.33. La pressione e la temperatura a 850~hPa si comportano allo stesso modo, con 1.38 e 1.58 a sette celle; i venti e l'umidità specifica stanno in mezzo.

\begin{figure}[htbp]
    \centering
    \phantomcaption
    \label{fig:ris-eventi-profilo-tipi}
    \begin{subfigure}[t]{0.437\textwidth}
        \centering
        \includegraphics[width=\linewidth]{era5_eventi_profilo_tipi_a}
        \caption[Percentile mediano di $G_{\mathrm{n}}$ secondo la distanza in celle dal bordo del cluster, per tipo di cluster, sugli anni di test: velocità verticale (rosa), precipitazione (viola), copertura nuvolosa (grigio scuro), temperatura a due metri (ocra), altri tipi in grigio chiaro.]{Percentile mediano di $G_{\mathrm{n}}$ secondo la distanza in celle dal bordo del cluster, per tipo di cluster, sugli anni di test: velocità verticale (rosa), precipitazione (viola), copertura nuvolosa (grigio scuro), temperatura a due metri (ocra), altri tipi in grigio chiaro. Le curve quasi coincidono.}
        \label{fig:ris-eventi-profilo-tipi-a}
    \end{subfigure}
    \hfill
    \begin{subfigure}[t]{0.523\textwidth}
        \centering
        \includegraphics[width=\linewidth]{era5_eventi_profilo_tipi_b}
        \caption[Scarto standardizzato mediano, in valore assoluto, della variabile del tipo, secondo la distanza in celle dal bordo del cluster, per ogni tipo di cluster, sugli anni di test.]{Scarto standardizzato mediano, in valore assoluto, della variabile del tipo, secondo la distanza in celle dal bordo del cluster, per ogni tipo di cluster, sugli anni di test. Alcune variabili tornano ai livelli di fondo entro due celle, altre restano a lungo lontane dal normale.}
        \label{fig:ris-eventi-profilo-tipi-b}
    \end{subfigure}
\end{figure}

\begin{figure}[htbp]
    \ContinuedFloat
    \centering
    \begin{subfigure}[t]{0.970\textwidth}
        \centering
        \includegraphics[width=\linewidth]{era5_eventi_profilo_tipi_c}
        \caption[Percentile mediano di ogni score a due celle fuori dal bordo dei propri cluster, un pallino per tipo di cluster, sugli anni di test; il segmento va dal tipo più basso al più alto e il numero a destra è la loro differenza.]{Percentile mediano di ogni score a due celle fuori dal bordo dei propri cluster, un pallino per tipo di cluster, sugli anni di test; il segmento va dal tipo più basso al più alto e il numero a destra è la loro differenza. $G_{\mathrm{n}}$ è fra i metodi più uniformi.}
        \label{fig:ris-eventi-profilo-tipi-c}
    \end{subfigure}
\end{figure}

Il risultato è coerente con il modo in cui queste variabili cambiano nello spazio, anche se la verifica non lo dimostra: precipitazione e nuvole possono cambiare da una cella all'altra, mentre temperatura e pressione variano su aree più ampie. Nel primo gruppo il nucleo di $G_{\mathrm{n}}$ contiene quasi tutta l'area in cui lo scarto standardizzato è alto. Nel secondo lo scarto standardizzato continua oltre il nucleo: $G_{\mathrm{n}}$ ne segnala solo una parte, e le celle vicine, pur lontane dal normale, restano sotto la sua soglia.
